对于 $m=1$，余子式矩阵恒等于一。其余情形令 $j_0=m-1$。由引理~\ref{ms:lemma-3-2}，$C$ 是次数 $j_0$ 的齐次多项式。令 $Q(M_1,\ldots,M_{j_0})$ 为 $t_1\cdots t_{j_0}$ 在 $C(\sum_jt_jM_j)$ 中的系数，从而将其极化。所得映射为多线性，且取值为 $j_0!C(M)$，其中输入取对角线 $M_1=\cdots=M_{j_0}=M$。

取 $M_j=g_jg_j^*$，其中 $g_j$ 为四元数列向量。若所有 $t_j>0$，令相应的非负二次型为零可知 $\sum_jt_jg_jg_j^*$ 的核是各 $g_j$ 张成空间的正交补。其像就是该张成空间，维数至多为 $m-1$。因此，其余子式矩阵作用于每个 $g_j$ 都为零，见引理~\ref{ms:lemma-3-2}。等式关于 $t$ 是多项式恒等式，因此系数 $Q(g_1g_1^*,\ldots,g_{j_0}g_{j_0}^*)$ 作用于每个 $g_j$ 亦为零。该系数为 Hermitian 矩阵，故还有 $g_j^*Q=0$。

以 $u_j(z)=\exp(\xi_j\cdot z)$ 为输入计算极化后的行散度，其中 $\xi_j$ 是任意实余向量，内积取实坐标点积。记 $(g_j)_a=\sum_\alpha e_\alpha(\xi_j)_{a\alpha}$。引理~\ref{ms:lemma-3-1} 中的 Hessian 符号给出
\[
Q(H(u_1),\ldots,H(u_{j_0}))
=e^{(\sum_j\xi_j)\cdot z}Q(g_1g_1^*,\ldots,g_{j_0}g_{j_0}^*).
\]
除去指数因子后，行散度为 $(\sum_jg_j)^*Q=0$。对于一般的光滑输入，该行散度是一个常系数多线性微分表达式。代入指数函数测试并除以各指数函数的乘积，可得到一个多项式，其中关于 $\xi_j$ 的各单项式编码对应输入的导数。它对一切实余向量都为零，故其所有系数均为零。令每个输入都等于 $u$，再除以 $j_0!$ 即得结论。
